\documentclass{article}
\usepackage{amsmath, amssymb, braket}

\title{Математическая модель GRA-Ethnos}
\author{GRA-Ethnos Team}
\date{\today}

\begin{document}
\maketitle

\section{Состояния и иерархия}

Пусть имеется множество агентов $\mathcal{A} = \{a_1,\dots,a_N\}$. Каждый агент $a$ описывается состоянием $\ket{\Psi_a^{(0)}} \in \mathcal{H}^{(0)}$, где $\mathcal{H}^{(0)}$ — гильбертово пространство индивидуальных состояний.

Коллективное состояние на уровне $l=1$ определяется как тензорное произведение:
\begin{equation}
\ket{\Psi^{(1)}} = \bigotimes_{a \in \mathcal{A}} \ket{\Psi_a^{(0)}} \in \mathcal{H}^{(1)} = \bigotimes_{a \in \mathcal{A}} \mathcal{H}^{(0)}.
\end{equation}

Аналогично вводятся состояния высших уровней (культура, идентичность) как тензорные произведения или более сложные конструкции.

\section{Цели и проекторы}

На каждом уровне $l$ задана цель $G_l$ — подпространство в $\mathcal{H}^{(l)}$, и проектор $\mathcal{P}_{G_l}$ на это подпространство. Условие согласованности иерархии:
\begin{equation}
\mathcal{P}_{G_l} = \bigotimes_{m<l} \mathcal{P}_{G_m}, \quad l\ge 1.
\end{equation}

\section{Пена}

Пена на уровне $l$ для состояния $\ket{\Psi^{(l)}}$ определяется как мера недиагональности в базисе, диагонализующем $\mathcal{P}_{G_l}$:
\begin{equation}
\Phi^{(l)}(\Psi^{(l)}, G_l) = \sum_{a \neq b} |\bra{\Psi_a^{(l)}} \mathcal{P}_{G_l} \ket{\Psi_b^{(l)}}|^2.
\end{equation}
Здесь $\ket{\Psi_a^{(l)}}$ — состояния подсистем на уровне $l$.

В частном случае уровня 1 (рой) это выражение упрощается до суммы квадратов недиагональных элементов матрицы резонанса между агентами.

\section{Функционал народа}

Общий функционал, минимизируемый в процессе эволюции:
\begin{equation}
J_{\text{ethnos}}(\Psi) = \sum_{l=0}^{3} \Lambda_l \Phi^{(l)}(\Psi^{(l)}, G_l) + \lambda_{\text{law}} \mathcal{L},
\end{equation}
где $\mathcal{L}$ — штраф за нарушение Аланского закона (например, доля запрещённых модификаций). Веса $\Lambda_l$ убывают с уровнем: $\Lambda_l = \lambda_0 \alpha^l$, $0<\alpha<1$.

\section{Теорема о субъектности}

\textbf{Теорема.} Если существует состояние $\Psi^*$, такое что $J_{\text{ethnos}}(\Psi^*) = 0$, то система $\mathcal{N} = (\mathcal{A}, \mathcal{G}, \mathcal{R}, \mathcal{L}, \mathcal{E})$ обладает свойством эмерджентной субъектности: она способна к целеполаганию, самосохранению и развитию, не сводимым к целям отдельных агентов.

\textit{Доказательство.} Условие $J_{\text{ethnos}}=0$ означает $\Phi^{(l)}=0$ для всех $l$, т.е. система находится в собственном состоянии всех проекторов. Это подразумевает внутреннюю согласованность на всех уровнях, включая идентичность. Следовательно, система может быть описана как единый субъект с волей к сохранению этого состояния. \hfill $\square$

\end{document}